\chapter{From machines to trees: UCT and the sign that flips}
\label{ch:uct}

\begin{objectives}
\begin{itemize}
\item Turn a single bandit problem into a tree of bandit problems, and state the four phases
  of a Monte-Carlo tree search iteration.
\item Perform backpropagation with correct sign flips, by hand, and explain why $Q$ in a
  Black-to-move node being $+0.9$ means Black is winning.
\item Explain why Leela evaluates a leaf with one network call instead of playing the game
  out at random, and what was wrong with the random version.
\item Compute why UCT without prior knowledge cannot function at chess's branching factor
  --- the argument that forces the next chapter.
\end{itemize}
\end{objectives}

\section{Bandits all the way down}

Chapter~\ref{ch:manymachines} solved one decision: given $K$ options with unknown values,
how do you spend $T$ pulls? A chess position is exactly that decision --- and so is every
position it leads to.

\begin{keyidea}
\textbf{A search tree is a bandit problem at every node.} At the root you must choose among
your legal moves. Whichever you choose, you arrive at a position where your \emph{opponent}
faces the same kind of choice, and so on. UCT (Kocsis \& Szepesv\'ari, 2006) is the
observation that you can run UCB1 at every node simultaneously, and that the whole thing
coheres: the value estimate at a node is fed by the estimates below it.
\end{keyidea}

The vocabulary, fixed once:

\begin{notationbox}
A \textbf{node} is a position the search has stored. An \textbf{edge} is a legal move out of
a stored position. Each edge carries $N(s,a)$ (visits) and $W(s,a)$ (accumulated value), with
$Q(s,a) = W(s,a)/N(s,a)$. A \textbf{leaf} is a node in the tree with no children yet. The
\textbf{root} is the position we must move in. $N(s) = \sum_b N(s,b)$ is the node's total
visits.
\end{notationbox}

\section{One iteration, four phases}

Every Monte-Carlo tree search --- Leela's included --- repeats the same four-phase cycle
until the budget runs out. One cycle adds exactly one new node to the tree.

\begin{center}
\begin{tikzpicture}[
   nd/.style={circle, draw, minimum size=7mm, inner sep=0pt, font=\footnotesize},
   newnd/.style={nd, fill=IdeaGreen!20, draw=IdeaGreen, thick},
   lbl/.style={font=\footnotesize\itshape, align=center},
   x=17mm, y=-14mm]
% tree
\node[nd, fill=BuildBlue!15] (r) at (0,0) {$s_0$};
\node[nd] (a) at (-1.1,1) {$s_1$};
\node[nd] (b) at (0.4,1) {};
\node[nd] (c) at (1.5,1) {};
\node[nd] (d) at (-1.7,2) {};
\node[nd] (e) at (-0.5,2) {$s_2$};
\node[newnd] (f) at (-0.1,3) {$s_3$};
\draw[very thick, PitfallRed] (r) -- (a);
\draw[very thick, PitfallRed] (a) -- (e);
\draw[very thick, IdeaGreen, dashed] (e) -- (f);
\draw (r) -- (b); \draw (r) -- (c); \draw (a) -- (d);
% annotations
\node[lbl, right=22mm of r, anchor=west] (t1)
  {\textbf{1. Selection.} From the root, repeatedly take the\\ highest-scoring edge
   ($Q+U$) until you reach a node\\ whose chosen move has no child yet.};
\node[lbl, below=6mm of t1.south west, anchor=north west] (t2)
  {\textbf{2. Expansion.} Create that child, $s_3$. Ask the\\ network for its policy over
   $s_3$'s legal moves.};
\node[lbl, below=4mm of t2.south west, anchor=north west] (t3)
  {\textbf{3. Evaluation.} The same network call returns\\ $V(s_3)$ --- one forward pass, no
   rollout.};
\node[lbl, below=4mm of t3.south west, anchor=north west] (t4)
  {\textbf{4. Backpropagation.} Walk back up the red path,\\ adding $V(s_3)$ to every $W$
   and $1$ to every $N$,\\ \emph{flipping the sign at each ply}.};
\end{tikzpicture}
\end{center}

The remarkable thing about this loop is how little it contains. There is no chess in it
whatsoever except the move generator. All the chess lives in the network call in phase 3;
the loop's only job is deciding \emph{which} position gets that call next. This is why
Chapter~\ref{ch:problem} insisted that \texttt{lc0.exe} has no chess knowledge --- it is
literally this loop, plus engineering.

\subsection{Phase 3, and what Monte-Carlo used to mean}

The name "Monte-Carlo" is a historical fossil, and it misleads people constantly.

In the original method (Go programs of the mid-2000s), phase 3 was a \textbf{rollout}: from
the new leaf, play random legal moves until the game ends, and score the result $+1$, $0$ or
$-1$. That is where the randomness --- and the casino name --- came from. Averaging thousands
of random games gave a shockingly usable estimate in Go.

It is a disaster in chess, for a reason worth stating precisely: \textbf{chess outcomes are
determined by short forcing sequences, and random play destroys them}. From a position where
you are a queen up, random play throws the queen back within a few moves; from a position
with mate in one, random play finds the mate one time in thirty. The rollout's answer is
dominated by noise that has nothing to do with the position.

\begin{established}
AlphaGo (2016) replaced the rollout with a learned \textbf{value network}, and AlphaZero
(2017) dropped rollouts entirely. Leela follows AlphaZero: phase 3 is one forward pass
returning $(P, V)$, and no game is ever played out. The statistics are still Monte-Carlo ---
we are averaging samples --- but the samples are network evaluations of concrete positions
chosen deliberately, not random games.
\end{established}

\begin{pitfall}
Because of this, the phrase "Leela plays out thousands of games in its head" --- common in
chess commentary and present in the earlier draft guide this book replaces --- is simply
false. Leela never plays a game out. At 800 nodes it has looked at 800 positions, evaluated
each once, and averaged. If your tool ever says "simulations", make sure the user cannot
infer rollouts from it.
\end{pitfall}

\section{Backpropagation and the sign flip}

Phase 4 is where the bugs live. A value is always \emph{from the point of view of the player
to move at that node} (Chapter~\ref{ch:problem}, negamax). So as the value walks up the tree,
it must alternate sign at every ply.

\begin{equation}
  \label{eq:backup}
  \text{for each edge } (s,a) \text{ on the path, from the leaf upwards:}\quad
  N(s,a) \mathrel{+}= 1, \qquad W(s,a) \mathrel{+}= v,\qquad v \leftarrow -v .
\end{equation}

\begin{worked}{backing up a value through three plies}
The search selected: root ($s_0$, White to move) $\to$ move $a$ $\to$ $s_1$ (Black to move)
$\to$ move $b$ $\to$ $s_2$ (White to move) $\to$ move $c$ $\to$ new leaf $s_3$ (Black to
move). The network evaluates $s_3$ at $V(s_3) = -0.6$ --- \emph{Black to move there, and
Black is worse}, since values are always for the side to move.

Now walk up, flipping each time:
\begin{center}
\begin{tabular}{@{}lllr@{}}
\toprule
Edge & Whose node & Value added to $W$ & Meaning \\
\midrule
$(s_2, c)$ & White to move at $s_2$ & $-(-0.6) = +0.6$ & good for White \\
$(s_1, b)$ & Black to move at $s_1$ & $-(+0.6) = -0.6$ & bad for Black \\
$(s_0, a)$ & White to move at $s_0$ & $-(-0.6) = +0.6$ & good for White \\
\bottomrule
\end{tabular}
\end{center}
Every $N$ on the path also increments by 1. Sanity check the top and bottom: a position that
is bad for Black-to-move must be good for White at the root, and indeed both White nodes
recorded $+0.6$. If you ever compute a chain where a single leaf makes both sides happy, you
have a sign bug.
\end{worked}

\begin{keyidea}
\textbf{$Q$ is always relative to the side to move at the parent node.} When you read
$\texttt{Q: 0.94}$ on an edge out of a node where Black is to move, it means the move is
excellent \emph{for Black}. This trips up everyone reading engine output for the first time,
and it is the single most common source of an evaluation displayed with the wrong sign in a
chess GUI.
\end{keyidea}

We can verify the convention on real output rather than take my word for it:

\begin{realdata}[title={Real engine output: the same position from both sides}]
The king-and-pawn position (White Ke6, pawn e5, Black Ke8) evaluated with White to move and
then with Black to move --- the identical board, one tempo apart:
\begin{center}
\begin{tabular}{@{}lrrl@{}}
\toprule
 & \texttt{WL} at root & best move & Stockfish d30 \\
\midrule
White to move & $+0.942$ & Kd6 & mate in 12 for White \\
Black to move & $-0.948$ & Kf8 & mate in 16 for White \\
\bottomrule
\end{tabular}
\end{center}
Both numbers say "White is winning". The first says it as $+0.94$ (White to move, White
happy); the second says it as $-0.95$ (Black to move, Black unhappy). The sign flipped
because the mover changed, not because anything on the board did.
\end{realdata}

\section{UCT: UCB1 at every node}

Put UCB1 (Equation~\eqref{eq:ucb1}) on each node and you have UCT. The selection rule at
node $s$ is
\begin{equation}
  \label{eq:uct}
  a^{*} \;=\; \argmax_{a}\left[\, Q(s,a) \;+\; c\sqrt{\frac{\ln N(s)}{N(s,a)}} \,\right],
\end{equation}
applied repeatedly from the root until you fall off the bottom of the tree, which is phase 1.
An unvisited edge has $N(s,a) = 0$ and therefore infinite bonus, so \emph{every legal move at
a node must be tried once before any is tried twice}.

That last sentence is the whole problem.

\subsection{Why UCT alone cannot play chess}

\begin{worked}{the cost of ``try everything once''}
Your corpus says $b \approx 31$ legal moves per position (Chapter~\ref{ch:problem}). UCT must
visit every child once before deepening anywhere. So:

\begin{center}
\begin{tabular}{@{}lrl@{}}
\toprule
To have looked once at every line of depth\dots & costs (nodes) & \\
\midrule
1 ply  & $31$      & one visit each to your own moves \\
2 ply  & $961$     & \\
3 ply  & $29{,}791$ & \\
4 ply  & $923{,}521$ & \\
\bottomrule
\end{tabular}
\end{center}

A search budget of 800 nodes --- the kind this project uses for a live board --- cannot
complete \emph{two plies}. It would spend its first 31 visits on the root's children,
including every move that hangs a queen, and then have 769 left to distribute over a tree
whose second layer alone needs 961.
\end{worked}

And the visits it does spend are spent stupidly. UCB1 has no way of knowing that
$1.\text{g4}$ deserves less attention than $1.\text{e4}$ \emph{before measuring it}. It must
pay a full visit to discover what any club player knows at a glance.

\begin{keyidea}
UCT's fatal assumption is that \textbf{all options start equal}. That is right for slot
machines --- they are anonymous --- and catastrophically wrong for chess moves, which come
with enormous amounts of prior evidence about which ones are worth examining. The fix is not
a better bonus. The fix is to \textbf{stop starting from ignorance}.
\end{keyidea}

\subsection{What the fix has to look like}

Before turning the page, notice how strong the constraints are on any repair:

\begin{enumerate}
\item It must let the search \emph{skip} moves entirely at low visit counts --- with 800
  nodes and 31 legal moves, most moves must receive zero visits, forever.
\item It must nevertheless be able to \emph{come back} to a skipped move if the favoured ones
  fail. A hard "only look at the top 3" cutoff cannot do this, and would make the engine
  blind exactly where blindness is fatal.
\item It must degrade gracefully into UCT-like behaviour as visits grow, so that with enough
  compute the prior's opinion can be overruled by measurement.
\end{enumerate}

Requirement 2 is the interesting one, and it is the reason the answer is not simply
"pre-filter the move list". Chapter~\ref{ch:puct} builds a bonus that satisfies all three ---
and the move that comes back from the dead in our real data is a queen sacrifice that the
network initially rated at $1.6\%$.

\begin{repobox}[title={in this repository}]
The four phases are the loop inside \texttt{engine/lc0.exe}; you never see them directly, but
every number in \texttt{--verbose-move-stats} is produced by them. When this project's
backend requests a search and reads back visit counts and principal variations, it is reading
the state of the tree that this loop left behind.
\end{repobox}

\begin{exercises}

\exhead[the loop]{ch5:loop} Write out the four phases in your own words, in one sentence
each, without looking. Then say which phase the neural network participates in, and how many
times it is called per iteration.

\exhead[sign flips]{ch5:signs} A selection path runs root (White) $\to s_1$ (Black)
$\to s_2$ (White) $\to s_3$ (Black) $\to$ new leaf $s_4$ (White), and the network reports
$V(s_4) = +0.35$. Write out the value added to $W$ on each of the four edges, with signs, and
say in words what the root edge's contribution means.

\exhead[reading a Q]{ch5:readq} An edge out of a Black-to-move node shows
\texttt{N: 210 (Q: 0.77)}. Is Black doing well or badly? A GUI displays this as "$+0.77$
White". What has gone wrong, and where in Equation~\eqref{eq:backup} did it go wrong?

\exhead[rollouts]{ch5:rollouts} Estimate, roughly, the probability that a random legal
continuation from a position with mate in one actually plays the mate, assuming 31 legal
moves. Then estimate the probability that a 40-ply random rollout preserves a decisive queen
advantage. Use these to argue in two sentences why rollouts died.

\exhead[UCT's cost]{ch5:uctcost} With $b = 31$, how many nodes does UCT need before it has
visited every 3-ply line once? Express the answer as a multiple of an 800-node budget. Now
redo it for the endgame branching factor $b = 15.8$ from your corpus, and comment.

\exhead[the first sweep]{ch5:firstsweep} Under pure UCT with 800 nodes at $b=31$, what
fraction of the budget is consumed just by giving every root move its first visit? If 20 of
those 31 moves are obviously bad to a club player, how much of the budget was wasted, and
what would you have to know in advance to avoid the waste?

\exhead[designing the repair]{ch5:designfix} Before reading Chapter~\ref{ch:puct}: propose a
modification of Equation~\eqref{eq:uct} that uses a prior probability $P(a\given s)$. Write
down your formula. Then check it against the three requirements in §5.5.2 and note which
ones it fails. (Keep this; you will compare it with the real answer.)

\exhead[the hard cutoff]{ch5:cutoff} Suppose you implement "only ever search the top 3 moves
by prior". Describe a concrete chess situation in which this loses a game outright, and
explain which of the three requirements it violates. Why is this failure mode particularly
dangerous for a \emph{training} tool as opposed to a playing engine?

\exhead[budget and depth]{ch5:budgetdepth} If a search concentrates $90\%$ of its visits on
one move at each ply, how deep can 800 nodes reach along the main line? Compare with the UCT
answer from Exercise~\ref{ex:ch5:uctcost}. This is a preview of the actual mechanism ---
state in one sentence what the search must be buying with that concentration.

\end{exercises}
